\section*{Chapitre \ref{ch:g8:winds-global-circulation} --- Les vents et la circulation générale}

\begin{solution}{exo:g8:winds-global-circulation:1}
(a) L'alizé du nord-est, qui souffle du nord-est. (b) Les \omterm{def:g8:winds-global-circulation:global-winds}{vents d'ouest},
qui soufflent du nord-ouest. (c) Les \omterm{def:g8:winds-global-circulation:global-winds}{vents d'est polaires}, qui soufflent
de l'est (du nord-est). (d) L'alizé du sud-est, qui souffle du sud-est.
\end{solution}

\begin{solution}{exo:g8:winds-global-circulation:2}
(a) L'air chaud monte et l'air froid descend. (b) Basse. (c) Des hautes
pressions vers les basses pressions. (d) Vers la droite.
\end{solution}

\begin{solution}{exo:g8:winds-global-circulation:3}
(a) Faux~: une \omterm{def:g8:winds-global-circulation:breezes}{brise de mer} souffle de la \omterm{def:g2:water-on-earth:ocean}{mer} vers la terre (le jour).
(b) Vrai. (c) Faux~: ils soufflent de l'est (du nord-est ou du sud-est).
(d) Vrai.
\end{solution}

\begin{solution}{exo:g8:winds-global-circulation:4}
Le dessin montre une boucle~: le long du sol, une flèche de la \omterm{def:g2:water-on-earth:ocean}{mer} vers la
terre~; au-dessus de la terre, une flèche qui monte~; en \omterm{def:g4:world-physical-map:altitude}{altitude}, une
flèche qui revient de la terre au-dessus de la \omterm{def:g2:water-on-earth:ocean}{mer}~; au-dessus de la \omterm{def:g2:water-on-earth:ocean}{mer},
une flèche qui descend. A est marqué au sol sous l'air qui descend
(au-dessus de la \omterm{def:g2:water-on-earth:ocean}{mer}), D sous l'air qui monte (au-dessus de la terre).
\end{solution}

\begin{solution}{exo:g8:winds-global-circulation:5}
(a) 1.41~fois plus de sol. (b) 50~\%. (c) 26~\%. (d) Près des \omterm{def:g3:continents-and-oceans:poles}{pôles}, la
lumière du Soleil arrive en biais et s'étale sur une grande surface~:
chaque mètre carré reçoit donc peu d'énergie.
\end{solution}

\begin{solution}{exo:g8:winds-global-circulation:6}
(a) Les \omterm{def:g8:winds-global-circulation:global-winds}{vents d'est polaires} (de \ang{70} à \ang{60} N), les \omterm{def:g8:winds-global-circulation:global-winds}{vents d'ouest}
(de \ang{60} à \ang{30} N), les \omterm{def:g8:winds-global-circulation:global-winds}{alizés} du nord-est (de \ang{30} N à
l'\omterm{def:g3:continents-and-oceans:equator}{équateur}), les \omterm{def:g8:winds-global-circulation:global-winds}{alizés} du sud-est (de l'\omterm{def:g3:continents-and-oceans:equator}{équateur} à \ang{30} S), les \omterm{def:g8:winds-global-circulation:global-winds}{vents
d'ouest} (de \ang{30} à \ang{60} S), les \omterm{def:g8:winds-global-circulation:global-winds}{vents d'est polaires} (de \ang{60}
à \ang{70} S). (b) Près de l'\omterm{def:g3:continents-and-oceans:equator}{équateur}, dans le pot-au-noir. (c) Vers
\ang{30} N et \ang{30} S, sous l'air qui descend.
\end{solution}

\begin{solution}{exo:g8:winds-global-circulation:7}
(a) La terre ($31 > 22$)~: le vent souffle de la \omterm{def:g2:water-on-earth:ocean}{mer} vers la terre, c'est
une \omterm{def:g8:winds-global-circulation:breezes}{brise de mer}. (b) La \omterm{def:g2:water-on-earth:ocean}{mer} ($21 > 17$)~: le vent souffle de la terre
vers la \omterm{def:g2:water-on-earth:ocean}{mer}, c'est une \omterm{def:g8:winds-global-circulation:breezes}{brise de terre}. (c) La terre~: $31 - 17 = 14$~degrés~;
la \omterm{def:g2:water-on-earth:ocean}{mer}~: $22 -
21 = 1$~degré. (d) Les rayons du Soleil ne chauffent qu'une mince couche de
sol, qui se réchauffe et se refroidit vite~; la \omterm{def:g2:water-on-earth:ocean}{mer} répartit sa chaleur
dans une épaisse couche d'eau, si bien que sa température change à peine
entre le jour et la nuit.
\end{solution}

\begin{solution}{exo:g8:winds-global-circulation:8}
(a) Du nord. (b) Vers la droite. (c) Du nord-est (c'est l'alizé du
nord-est). (d) Sans la rotation, il soufflerait du sud~; la rotation le
dévie vers la gauche~; il souffle du sud-est.
\end{solution}

\begin{solution}{exo:g8:winds-global-circulation:9}
Près de l'\omterm{def:g3:continents-and-oceans:equator}{équateur}, le sol et la \omterm{def:g2:water-on-earth:ocean}{mer} sont fortement \emph{chauffés} par le
Soleil~; l'air chaud et humide \emph{monte}~; en montant, il \emph{se refroidit},
et sa vapeur d'eau se change en \emph{\omterm{def:g4:clouds-rain-snow:cloud}{nuages}}, qui donnent des averses et
des orages. Vers \ang{25} N, l'air qui est monté à l'\omterm{def:g3:continents-and-oceans:equator}{équateur}
\emph{descend}~; sous lui, la pression au sol est \emph{haute}~; l'air qui
descend se réchauffe, ses \omterm{def:g4:clouds-rain-snow:cloud}{nuages} se dissipent et la région reste
\emph{sèche}~: un \omterm{def:g5:weather-and-climate:desert}{désert}.
\end{solution}

\begin{solution}{exo:g8:winds-global-circulation:10}
(a) Les \omterm{def:g8:winds-global-circulation:global-winds}{vents d'ouest}. Non~: ils soufflent de l'ouest, contre un navire
qui va vers l'ouest. (b) Pour atteindre la bande des \omterm{def:g8:winds-global-circulation:global-winds}{alizés}, au sud
d'environ \ang{30} N. (c) L'alizé du nord-est, entre l'\omterm{def:g3:continents-and-oceans:equator}{équateur} et environ
\ang{30} N. (d) Les \omterm{def:g8:winds-global-circulation:global-winds}{vents d'ouest}, au nord d'environ \ang{30} N (vers
\ang{40} N, la \omterm{def:g6:latitude-longitude:latitude}{latitude} du port d'attache).
\end{solution}

\begin{solution}{exo:g8:winds-global-circulation:11}
(a) $1 \div 0.5 = 2$~: le tableau donne 2. (b) $1 \div 0.866 \approx
1.1547$, soit \textbf{1.15}. (c) À \ang{30}~: $1000 \div 1.15 \approx
\qty{870}{W}$ sur chaque mètre carré (exactement $1000 \times 0.866 =
\qty{866}{W}$). À \ang{60}~: $1000 \div 2 = \qty{500}{W}$.
\end{solution}

\begin{solution}{exo:g8:winds-global-circulation:12}
(a) Trois boucles~: l'air monte à l'\omterm{def:g3:continents-and-oceans:equator}{équateur} et descend vers \ang{30}~; il
monte vers \ang{60} et descend vers \ang{30}~; il monte vers \ang{60} et
descend au \omterm{def:g3:continents-and-oceans:poles}{pôle}. (b) D à l'\omterm{def:g3:continents-and-oceans:equator}{équateur} et vers \ang{60}~; A vers \ang{30} et
au \omterm{def:g3:continents-and-oceans:poles}{pôle}. (c) De \ang{30} à l'\omterm{def:g3:continents-and-oceans:equator}{équateur}, l'air de surface va vers le sud et
il est dévié vers la droite~: c'est l'alizé du nord-est. De \ang{30} à
\ang{60}, il va vers le nord et il est dévié vers la droite~: il vient du
sud-ouest. Du \omterm{def:g3:continents-and-oceans:poles}{pôle} à \ang{60}, il va vers le sud, dévié vers la droite~:
c'est le vent d'est polaire. (d) Dans la cellule du milieu, l'air de
surface s'écoule vers le \omterm{def:g3:continents-and-oceans:poles}{pôle}, c'est-à-dire vers le nord~; dévié vers la
droite, il se dirige vers le nord-est~: ce sont les \omterm{def:g8:winds-global-circulation:global-winds}{vents d'ouest}.
\end{solution}

\begin{solution}{pb:g8:winds-global-circulation:1}
\textbf{1.} $59.6 - 23.5 = 36.1$~degrés de \omterm{def:g6:latitude-longitude:longitude}{longitude}.

\textbf{2.} $36.1 \times 108 \approx 3899$~: environ \qty{3900}{km}.

\textbf{3.} Oui~: $3900 - 3896 = 4$~km, une toute petite différence.

\textbf{4.} Oui~: tous deux se trouvent vers \ang{14} N, entre l'\omterm{def:g3:continents-and-oceans:equator}{équateur}
et environ \ang{30} N, la bande de l'alizé du nord-est.

\textbf{5.} Du nord-est (les flèches sont tournées vers le sud-ouest).

\textbf{6.} Vers l'avant, et un peu sur le côté~: le vent vient de
l'arrière du navire, légèrement de sa droite, et le pousse vers l'ouest.

\textbf{7.} Il fait partie de la \omterm{def:g8:winds-global-circulation:circulation}{circulation générale}~: l'air ne cesse de
descendre vers \ang{30} N et de monter près de l'\omterm{def:g3:continents-and-oceans:equator}{équateur}~; l'air de
surface s'écoule donc sans cesse vers l'\omterm{def:g3:continents-and-oceans:equator}{équateur}, et la rotation de la
Terre le dévie toujours vers la droite.

\textbf{8.} L'alizé soufflerait contre le navire. Le capitaine doit
remonter vers le nord, au-delà d'environ \ang{30} N, pour trouver les
\omterm{def:g8:winds-global-circulation:global-winds}{vents d'ouest}.

\textbf{9.} $7 \times 24 = \qty{168}{km}$ par jour.

\textbf{10.} $3 \times 168 = \qty{504}{km}$.

\textbf{11.} $3896 \div 168 \approx 23.2$~jours.

\textbf{12.} Environ \textbf{23~jours}~: moins d'un mois.
\end{solution}
